\clearpage
\VerdierPhysicalPageAnchor{0119}%
\begingroup
\small
\setlength{\abovedisplayskip}{4pt plus 2pt minus 2pt}%
\setlength{\belowdisplayskip}{4pt plus 2pt minus 2pt}%

\VerdierSemanticAnchor{ch2-1-2-8-splitting-complete}
$fgf=f$ and $gfg=g$. A morphism is called \emph{splittable} if it has a
splitting. Decomposable morphisms are splittable. The converse is not
necessarily true (example: the category of free modules over a nonprincipal
Dedekind ring).

\VerdierSemanticAnchor{ch2-1-2-8-decomposable-category}
An additive category is called \emph{decomposable} if every morphism in that
category is decomposable. A decomposable category is abelian.

\VerdierSemanticAnchor{ch2-1-2-8-triangle-decomposition}
Let $X\xrightarrow{f}Y$ be a morphism in a triangulated category
$\mathcal D$, and let

\VerdierSemanticAnchor{ch2-1-2-8-decomposition-diagram-repeat}
\begin{equation*}
  \begin{tikzcd}[ampersand replacement=\&,column sep=5.2em,row sep=3.4em]
    X \arrow[r,"f"] \arrow[d,"\wr"']
      \& Y \arrow[d,"\wr"] \\
    Z\oplus W
      \arrow[r,"{\left(\begin{smallmatrix}0&\operatorname{id}_W\\0&0\end{smallmatrix}\right)}"]
      \& W\oplus Z'
  \end{tikzcd}\,.
\end{equation*}
be a decomposition of $f$. Every distinguished triangle containing $f$ is
isomorphic to the triangle

\VerdierSemanticAnchor{ch2-1-2-8-decomposed-triangle}
\begin{equation*}
  \begin{tikzcd}[ampersand replacement=\&,column sep=2.0em]
    Z\oplus W
      \arrow[r,"{\left(\begin{smallmatrix}0&\operatorname{id}_W\\0&0\end{smallmatrix}\right)}"]
      \& W\oplus Z'
      \arrow[r,"{\left(\begin{smallmatrix}0&\operatorname{id}_{Z'}\\0&0\end{smallmatrix}\right)}"]
      \& Z'\oplus Z[1]
      \arrow[r,"{\left(\begin{smallmatrix}0&\operatorname{id}_{Z[1]}\\0&0\end{smallmatrix}\right)}"]
      \& Z[1]\oplus W[1]
  \end{tikzcd}\,.
\end{equation*}

\VerdierSemanticAnchor{ch2-1-2-8-decomposed-triangle-derivation}
This follows immediately from (1.2.4), (1.2.5), and axioms TRI and TRII.

\VerdierSemanticAnchor{ch2-1-2-9}
\medskip
\noindent\textbf{Proposition 1.2.9.}\quad
\emph{Let $\mathcal D$ be a triangulated category and let
$f\colon X\longrightarrow Y$ be a morphism in $\mathcal D$. Consider the
following properties of the morphism $f$:}

\VerdierSemanticAnchor{ch2-1-2-9-properties}
\begin{enumerate}
\renewcommand{\labelenumi}{\roman{enumi})}
\item \emph{The morphism $f$ has a kernel.}
\item \emph{The morphism $f$ has a cokernel.}
\item \emph{The morphism $f$ is decomposable.}
\item \emph{The morphism $f$ is splittable.}
\end{enumerate}

\noindent\emph{Then the following implications hold:}

\VerdierSemanticAnchor{ch2-1-2-9-implications}
\[
  \mathrm{i)}\Longleftrightarrow\mathrm{ii)}
  \Longleftrightarrow\mathrm{iii)}\Longrightarrow\mathrm{iv)}\,.
\]

\VerdierSemanticAnchor{ch2-1-2-9-countable-equivalence}
\emph{If countable products or countable sums are representable in
$\mathcal D$, properties \textnormal{i)}, \textnormal{ii)},
\textnormal{iii)}, and \textnormal{iv)} are equivalent.}

\VerdierSemanticAnchor{ch2-1-2-9-proof-opening}
First, we plainly have the implications
$\mathrm{iii)}\Longrightarrow\mathrm{i)}$,
$\mathrm{iii)}\Longrightarrow\mathrm{ii)}$,
$\mathrm{iii)}\Longrightarrow\mathrm{iv)}$. The implication
$\mathrm{i)}\Longrightarrow\mathrm{iii)}$ is equivalent to the implication
\endgroup
